\documentclass{ximera}

\input{../../preamble.tex}

\title{Core Analysis}

\begin{document}

\begin{abstract}
Part 2
\end{abstract}
\maketitle




As a starting point for function analysis, let's just list characteristics of our Core Functions. A full analysis would also provide reasoning.


\begin{explanation} \textbf{\textcolor{red!75!green}{Exponential Functions}} 


\[ E(x) = r^x  \, \text{ where } \, r > 0  \, \text{ and } \, r \ne 1 \]



\textbf{Domain:} The domain $(-\infty, \infty)$.

\textbf{Zeros:} Core Exponential Functions have no zeros.



\textbf{Continuity:} Core Exponential Functions are continuous.



\textbf{End-Behavior:} 



If $r > 1$, the Core Exponential Functions tend to $0$ as the domain becomes unbounded negatively. Core Exponential Functions tend to $\infty$ as the domain becomes unbounded positively.



\[
\lim\limits_{x \to \infty} r^x = \infty
\]


\[
\lim\limits_{x \to -\infty} r^x = 0
\]


If $0 < r < 1$, then this is reversed.



\textbf{Behavior (Inc/Dec):} 


If $r > 1$, then Core Exponential Functions are increasing functions.

If $0 < r < 1$, then Core Exponential Functions are decreasing functions.





\textbf{Global Maximum and Minimum:} Core Exponential Functions have no global maximum or minimum.

\textbf{Local Maximum and Minimum:}  Core Exponential Functions have no local maximums or minimums.

\textbf{Range:} The range of a Core Exponential Function is either $(0,\infty)$ or $(-\infty, 0)$.

\end{explanation}











\begin{explanation} \textbf{\textcolor{red!75!green}{Core Logarithmic Functions}} 


\[ L(x) = \log_r(x)  \]



\textbf{Domain:} The domain $(0, \infty)$.

\textbf{Zeros:} Core Logarithmic Function have $1$ as their only zero.



\textbf{Continuity:} Core Logarithmic Functions are continuous.



\textbf{End-Behavior:} Since the domain is $(0, \infty)$, Core Logarithmic Functions only end-behavior as the domain become unbounded positively.


If $r > 1$, then

\[
\lim\limits_{x \to \infty} \log_r(x) = \infty
\]


The other side is called Singularity Behavior.

\[
\lim\limits_{x \to 0^+} \log_r(x) = -\infty
\]





If $0 < r < 1$, then

\[
\lim\limits_{x \to \infty} \log_r(x) = -\infty
\]


The other side is called Singularity Behavior.

\[
\lim\limits_{x \to 0^+} \log_r(x) = \infty
\]







\textbf{Behavior (Inc/Dec):} 




If $r > 1$, then Core Logarithmic Functions are increasing functions.


If $0 < r < 1$, then Core Logarithmic Functions are decreasing functions.





\textbf{Global Maximum and Minimum:} Core Logarithmic Functions have no global maximum or minimum.

\textbf{Local Maximum and Minimum:}  Core Logarithmic Functions have no local maximums or minimums.

\textbf{Range:} The range of Core Logarithmic Functions is $(-\infty,\infty)$.

\end{explanation}

















\begin{explanation} \textbf{\textcolor{red!75!green}{Sine and Cosine}} 


\[ S(x) = \sin(x)  \]

\[ C(x) = \cos(x)  \]



\textbf{Domain:} The domain both functions is $(-\infty, \infty)$.

\textbf{Zeros:} Both functions have zeros and since both functions are periodic, they have an infinite number of zeros.


\begin{itemize}
\item Zeros of $\sin(x)$ are ${k \pi \, | \, k \in \mathbb{Z}}$.
\item Zeros of $\cos(x)$ are ${\frac{\pi}{2} + k \pi\, | \, k \in \mathbb{Z}}$.
\end{itemize}


\textbf{Continuity:} Both functions are continuous.



\textbf{End-Behavior:} Both functions continue to oscillate.






\textbf{Behavior (Inc/Dec):} Both functions switch back-and-forth between increasing and decreasing with a period of $2\pi$.

\begin{itemize}
\item $\sin(x)$ increases on $\left(0, \frac{\pi}{2} \right)$
\item $\sin(x)$ decreases on $\left(\frac{\pi}{2}, \frac{3\pi}{2} \right)$
\item $\sin(x)$ increases on $\left(\frac{3\pi}{2}, 2\pi \right)$
\end{itemize}

\begin{itemize}
\item $\cos(x)$ decreases on $(0, \pi)$
\item $\cos(x)$ increases on $(\pi. 2\pi)$
\end{itemize}







\textbf{Global Maximum and Minimum:} Both functions have a global maximum of $1$ and a global minimum of $-1$.




\textbf{Local Maximum and Minimum:}  Both functions $1$ and $-1$ as their local extrema.




\textbf{Range:} The range of both functions is $[-1, 1]$.

\end{explanation}



















\begin{explanation} \textbf{\textcolor{red!75!green}{Absolute Value Function}} 


\[ Abs(x) = |x|  \]



\textbf{Domain:} The domain is $(-\infty, \infty)$.

\textbf{Zeros:} $0$ is the only zero.




\textbf{Continuity:} The Core Absolute Value Function is continuous.

\textbf{End-Behavior:} The Core Absolute Value Function becomes unbounded positively on both sides.






\textbf{Behavior (Inc/Dec):} 

The Core Absolute Value Function decreases on $(-\infty, 0)$ and increases on $(0, \infty)$







\textbf{Global Maximum and Minimum:} The Core Absolute Value Function has $0$ as its global minimum and it has no global maximum.




\textbf{Local Maximum and Minimum:}  The Core Absolute Value Function has $0$ as its only local minimum and it has no local maximum.



\textbf{Range:} The range is $[0, \infty)$.

\end{explanation}






































\begin{center}
\textbf{\textcolor{green!50!black}{ooooo-=-=-=-ooOoo-=-=-=-ooooo}} \\

more examples can be found by following this link\\ \link[More Examples of Visual Behavior]{https://ximera.osu.edu/csccmathematics/precalculus1/precalculus1/visualBehavior/examples/exampleList}

\end{center}





\end{document}
